\documentclass[12pt,letterpaper,onecolumn]{exam}
\usepackage[utf8]{inputenc}
\usepackage[T1]{fontenc}
\usepackage[brazil]{babel}
\usepackage{amsmath}
\usepackage{amssymb}
\usepackage[lmargin=71pt,tmargin=1.2in]{geometry}
\usepackage{enumitem}
\usepackage{booktabs}

\thispagestyle{empty}

\begin{document}

\begingroup
    \centering
    \LARGE Exercícios Extras - Econometria I - 2026\\[0.5em]
    \large Selecionados das Listas 5 a 11 \par
\endgroup
\rule{\textwidth}{0.4pt}

\begin{questions}

\question \textbf{L10 (2(8.7))} Considere um modelo em nível de empregado,
\[
y_{i,e} = \beta_0 + \beta_1 x_{i,e,1} + \beta_2 x_{i,e,2} + \cdots + \beta_k x_{i,e,k} + f_i + v_{i,e},
\]
em que a variável não observada \(f_i\) é um ``efeito da empresa'' para cada empregado da empresa \(i\). O termo de erro \(v_{i,e}\) é específico do empregado e da empresa. O erro de combinação é \(u_{i,e} = f_i + v_{i,e}\), como na equação (8.28).
\begin{enumerate}[label=(\roman*)]
    \item Suponha que \(\operatorname{Var}(f_i)=\sigma_f^2\), \(\operatorname{Var}(v_{i,e})=\sigma_v^2\), e que \(f_i\) e \(v_{i,e}\) sejam não correlacionadas. Demonstre que \(\operatorname{Var}(u_{i,e})=\sigma_f^2+\sigma_v^2\). Chame esse valor de \(\sigma^2\).
    \item Agora suponha que, para \(e \neq g\), \(v_{i,e}\) e \(v_{i,g}\) sejam não correlacionadas. Demonstre que \(\operatorname{Cov}(u_{i,e},u_{i,g})=\sigma_f^2\).
    \item Seja
    \[
    \bar{u}_i = m_i^{-1}\sum_{e=1}^{m_i} u_{i,e}
    \]
    a média dos erros de combinação em determinada empresa. Demonstre que
    \[
    \operatorname{Var}(\bar{u}_i)=\sigma_f^2+\frac{\sigma_v^2}{m_i}.
    \]
    \item Detalhe a relevância do item (iii) para a estimação por MQP usando dados nivelados pela média da empresa, em que o peso usado para a observação \(i\) é o tamanho da empresa.
\end{enumerate}

\vspace{0.7cm}

\question \textbf{L6 (2(ANPEC, 2012))} Usando uma base de dados que tem informações de 65.535 trabalhadores, queremos verificar se existe desigualdade salarial entre os setores da economia. Considere que a economia está dividida em quatro setores: indústria, comércio, serviços e construção. Cada um dos trabalhadores está em um dos quatro setores e eles são mutuamente exclusivos. Seja \(Y_i\) o salário mensal do trabalhador \(i\) e definimos para cada setor uma variável binária que é igual a 1 se o trabalhador pertence ao setor e 0 caso contrário. Estimando um modelo linear de regressão, obtemos o seguinte resultado:

\[
\hat{Y}_i = 4{,}00 + 0{,}12\,educ_i + 0{,}03\,idade_i + 0{,}40\,homem_i - 0{,}05\,DI_i - 0{,}15\,DC_i - 0{,}25\,DCons_i
\]
\vspace{-0.8cm}
\[
\quad (0{,}02)\quad (0{,}008)\quad (0{,}0001)\qquad (0{,}0005)\qquad (0{,}0011)\qquad (0{,}003)\qquad (0{,}005)
\]
\[
R^2 = 0{,}83
\]

Em que \textit{educ} representa o número de anos de estudo de cada trabalhador, \textit{idade} é medida em anos, \textit{homem} é uma variável binária que assume valor igual a 1 se é homem e 0 caso contrário, \(DI\) representa a dummy para indústria, \(DC\) para comércio e \(DCons\) para construção. Entre parênteses encontra-se o erro-padrão de cada estimativa.

Baseado nas informações anteriores, julgue as seguintes afirmativas:
\begin{enumerate}[label=(\arabic*)]
    \item Com base nos resultados anteriores é possível rejeitar ao nível de 5\% de significância a hipótese nula de que o salário do setor da indústria é igual ao salário do setor de serviços para trabalhadores com o mesmo nível educacional, a mesma idade e do mesmo sexo. A hipótese alternativa é que os salários nestes setores sejam diferentes.
    \item Com base nos resultados anteriores é possível rejeitar ao nível de 5\% de significância a hipótese nula de que o salário no setor de construção é igual ao salário no setor de comércio, mantendo educação, idade e sexo fixos. A hipótese alternativa é que os salários nesses setores sejam diferentes.
    \item Com base nos resultados anteriores, é possível rejeitar ao nível de 5\% de significância a hipótese nula de que os salários nos 4 setores da economia são iguais, mantendo constante educação, idade e sexo.
    \item Os resultados do modelo anterior permitem testar a hipótese nula de que o retorno salarial entre homem e mulher é diferente para cada nível educacional, ao nível de 5\% de significância.
    \item Com base nos resultados anteriores, podemos testar a hipótese de que o intercepto do modelo linear de salário em função de educação, idade e setor para homem é diferente do intercepto do mesmo modelo linear de salário para mulher.
\end{enumerate}


\vspace{0.7cm}

\question \textbf{L9 (2(7.9))} Seja \(d\) uma variável \textit{dummy} (binária) e \(z\) uma variável quantitativa. Considere o modelo
\[
y = \beta_0 + \delta_0 d + \beta_1 z + \delta_1 d \cdot z + u;
\]
essa é uma versão geral de um modelo com uma interação entre uma variável \textit{dummy} e uma variável quantitativa. [Um exemplo está na equação (7.17).]
\begin{enumerate}[label=(\roman*)]
    \item Como isso não alterará nada importante, defina o erro com valor zero, \(u=0\). Então, quando \(d=0\), podemos escrever o relacionamento entre \(y\) e \(z\) como a função \(f_0(z)=\beta_0+\beta_1 z\). Escreva essa mesma relação quando \(d=1\), em que você deve usar \(f_1(z)\) no lado esquerdo para denotar a função linear de \(z\).
    \item Considerando \(\delta_1 \neq 0\) (o que significa que as duas não são paralelas), demonstre que o valor de \(z^\ast\) é tal que \(f_0(z^\ast)=f_1(z^\ast)\) e \(z^\ast = -\delta_0/\delta_1\). Esse é o ponto no qual as duas linhas se cruzam [como na Figura 7.2(b)]. Demonstre que \(z^\ast\) será positivo se, e somente se, \(\delta_0\) e \(\delta_1\) tiverem sinais opostos.
    \item Usando os dados contidos no arquivo \texttt{TWOYEAR}, a seguinte equação poderá ser estimada:
    \[
    \widehat{\log(wage)} = 2{,}289 - 0{,}357\,female + 0{,}050\,totcoll + 0{,}030\,female \cdot totcoll
    \]
    \[
    \hspace{2.0cm} (0{,}011)\hspace{1.2cm}(0{,}015)\hspace{1.1cm}(0{,}003)\hspace{1.0cm}(0{,}005)
    \]
    \[
    n = 6{,}763,\; R^2 = 0{,}202
    \]
    em que todos os coeficientes e erros-padrão foram arredondados com três casas decimais. Usando essa equação, encontre o valor de \textit{totcoll} de tal forma que os valores previstos de \(\log(wage)\) sejam os mesmos tanto para homens quanto para mulheres.
    \item Com base na equação do item (iii), as mulheres podem, de forma realística, obter educação superior de modo que seus ganhos alcancem os dos homens? Explique.
\end{enumerate}

\vspace{0.7cm}

\question \textbf{Questão nova: desempenho em matemática e especificação do modelo}

Suponha que uma pesquisadora esteja interessada em estudar o desempenho em matemática de escolas usando dados do arquivo \texttt{MEAPSINGLE}, de Wooldridge. A variável dependente é \textit{math4}, a porcentagem de alunos da quarta série aprovados no teste de matemática. A principal variável explicativa é \textit{lexppp}, o logaritmo dos gastos por estudante. A pesquisadora estima diferentes especificações por MQO. Os resultados hipotéticos estão na tabela abaixo.

\begin{center}
\small
\begin{tabular}{lccccc}
\toprule
& \multicolumn{5}{c}{Variável dependente: \textit{math4}} \\
\cmidrule(lr){2-6}
Regressores & (1) & (2) & (3) & (4) & (5) \\
\midrule
\textit{lexppp} & 6{,}85 & 9{,}01 & 1{,}93 & 68{,}40 & 11{,}40 \\
& (3{,}10) & (4{,}04) & (2{,}82) & (41{,}20) & (4{,}35) \\[0.15cm]
\textit{lexppp}\(^2\) & -- & -- & -- & -3{,}45 & -- \\
& -- & -- & -- & (2{,}18) & -- \\[0.15cm]
\textit{free} & -- & -0{,}422 & -0{,}060 & -0{,}419 & 0{,}820 \\
& -- & (0{,}071) & (0{,}054) & (0{,}071) & (0{,}410) \\[0.15cm]
\textit{lmedinc} & -- & -0{,}752 & -10{,}78 & -0{,}710 & -1{,}10 \\
& -- & (5{,}358) & (3{,}76) & (5{,}30) & (5{,}22) \\[0.15cm]
\textit{pctsgle} & -- & -0{,}274 & -0{,}397 & -0{,}270 & -0{,}250 \\
& -- & (0{,}161) & (0{,}111) & (0{,}160) & (0{,}158) \\[0.15cm]
\textit{read4} & -- & -- & 0{,}667 & -- & -- \\
& -- & -- & (0{,}042) & -- & -- \\[0.15cm]
\textit{lexppp}\(\times\)\textit{free} & -- & -- & -- & -- & -0{,}145 \\
& -- & -- & -- & -- & (0{,}052) \\[0.15cm]
Constante & 18{,}30 & 24{,}49 & 149{,}38 & -271{,}20 & 4{,}80 \\
& (25{,}10) & (59{,}24) & (41{,}70) & (193{,}60) & (61{,}50) \\
\midrule
Observações & 229 & 229 & 229 & 229 & 229 \\
\(R^2\) & 0{,}031 & 0{,}472 & 0{,}749 & 0{,}478 & 0{,}486 \\
\bottomrule
\end{tabular}
\end{center}

\noindent \textit{Variáveis:} \textit{math4} = porcentagem de aprovação no teste de matemática; \textit{lexppp} = log dos gastos por estudante; \textit{free} = porcentagem de alunos elegíveis para merenda gratuita; \textit{lmedinc} = log da renda mediana do distrito; \textit{pctsgle} = porcentagem de alunos de famílias monoparentais; \textit{read4} = porcentagem de aprovação no teste de leitura.

\begin{enumerate}[label=(\alph*)]
    \item Compare as estimativas do coeficiente de \textit{lexppp} nas colunas (1), (2) e (3). O que muda quando controles socioeconômicos são adicionados? O que muda quando \textit{read4} é adicionada?
    \item Usando a coluna (2), interprete o coeficiente de \textit{lexppp}. Qual é o efeito estimado de um aumento de 10\% nos gastos por estudante sobre \textit{math4}? Esse efeito é estatisticamente significante ao nível de 5\%?
    \item Compare as colunas (2) e (4). A inclusão do termo quadrático \textit{lexppp}\(^2\) parece importante? Com base na coluna (4), em que valor de \textit{lexppp} o efeito marginal dos gastos se torna zero?
    \item A coluna (3) tem \(R^2\) muito maior do que a coluna (2). Explique por que, apesar disso, a coluna (2) pode ser mais adequada caso o objetivo seja estimar o efeito causal dos gastos por estudante sobre o desempenho em matemática.
    \item Na coluna (5), o modelo inclui a interação \textit{lexppp}\(\times\)\textit{free}. Escreva o efeito marginal de \textit{lexppp} sobre \textit{math4}. Calcule esse efeito quando \textit{free}=20 e quando \textit{free}=60. O que essa especificação sugere sobre a relação entre gastos e desempenho em escolas com maior proporção de alunos pobres?
\end{enumerate}

\end{questions}

\end{document}
